% =====================================================================
\chapter{Conclusion, Limitation and Future Scope}
\label{ch:conclusion}
% =====================================================================

This chapter summarises what the project achieved, states the limitations of
the system and of the study, and describes the work that would most improve
\veriscript{}.

% ---------------------------------------------------------------------
\section{Conclusion}
% ---------------------------------------------------------------------

The project set out to build an open detector of AI-generated text that is
simultaneously accurate, inexpensive, explainable, robust and fair, by routing
each document through a cheap interpretable heuristic stage and escalating
only the ambiguous ones to a fine-tuned transformer. All five objectives of
Chapter~\ref{ch:problem} were carried out:
\begin{itemize}
  \item \textbf{O1, heuristic baseline.} A detector over 38 features (GPT-2
        token statistics, burstiness, stylometry and POS bigrams) with an
        XGBoost head and isotonic calibration reaches F1 0.946 and AUROC 0.988
        on the test split, with a well-calibrated probability (ECE 0.015).
  \item \textbf{O2, transformer baseline.} RoBERTa-base fine-tuned on a
        balanced corpus of 28,006 documents reaches F1 0.975 and AUROC 0.999
        in distribution.
  \item \textbf{O3, hybrid architecture.} A confidence gate with a band tuned
        on validation data, offered as three operating points.
  \item \textbf{O4, unified evaluation.} The three systems were compared on
        identical data for accuracy with bootstrap confidence intervals,
        calibration, a held-out generator family, leave-one-generator-out
        retraining, out-of-distribution and adversarial test sets, fairness
        and CPU and GPU latency.
  \item \textbf{O5, packaging.} The system is available as a Python library, a
        command line tool, a REST service, a web interface and a Docker image,
        with a test suite of 34 tests.
\end{itemize}

The evaluation leads to four conclusions.

\textbf{First, the two stages are complementary, and the cascade exploits
this.} With only 8\% of documents escalated, the thorough cascade matches the
transformer's F1 within statistical uncertainty (0.971 against 0.975) and makes
half as many false accusations as either baseline (49 against 94 and 100 out of
2,000 human texts). The heuristic is reliable when it is confident (97.7\% of
its fast-path decisions are correct), and the transformer is reliable exactly
where the heuristic is not. On validation data, cascades that escalate 20--30\%
of documents are more accurate than the transformer alone.

\textbf{Second, the fine-tuned transformer is far less robust than its
in-distribution scores suggest.} On out-of-distribution text (MAGE) it flags
69\% of human texts, and on student essays it flags more than half of the native
and nearly two thirds of the non-native writers. The heuristic, whose features
describe generic statistical properties of text, degrades much more gracefully.
A cascade that keeps confident heuristic decisions is therefore not only
cheaper in principle but considerably safer in practice than a stand-alone
transformer, and the gate visibly escalates more documents (up to 44\%) when it
meets unfamiliar data.

\textbf{Third, the efficiency premise does not hold in the implemented form.}
The heuristic stage relies on a GPT-2 pass that costs as much as the RoBERTa
classifier itself, plus a vocabulary projection and spaCy tagging, so the
cascade is about twice as slow as the transformer on the CPU and seven times as
slow on the GPU. The design principle ``cheap by default'' is only satisfied if
the first stage is genuinely cheaper than the second; a reference language
model of the same size as the classifier violates it. This is a negative result,
but a useful one: it identifies precisely which component must change
(Section~\ref{sec:future}).

\textbf{Fourth, fairness must be measured in absolute as well as relative
terms.} All systems flag non-native writing more often than native writing. A
disparity-only constraint was satisfied by the tuning procedure in a degenerate
way, by making the detector equally bad for both groups; the worst-group
constraint adopted instead protects non-native writers relative to the
transformer (25.5\% against 63.8\% false positives) but not relative to the
heuristic alone (17.0\%). No configuration is fair enough to justify automatic
sanctions.

In summary, \veriscript{} delivers an open, explainable and comparatively robust
detector that equals a fine-tuned transformer in accuracy while halving false
positives, and it provides a complete, reproducible evaluation framework. It
does not yet deliver the intended reduction in computational cost.

% ---------------------------------------------------------------------
\section{Limitation}
% ---------------------------------------------------------------------

\textbf{Limitations of the system.}
\begin{itemize}
  \item \textbf{Cost.} The heuristic stage is slower than the transformer it is
        meant to save (Table~\ref{tab:latency}); the cascade is therefore not
        suitable where throughput is the main concern.
  \item \textbf{False positives on unfamiliar genres.} Even the best
        configuration flags about one in four TOEFL essays and one in five
        human MAGE texts. The output must be treated as a prompt for human
        review and never as evidence of misconduct on its own.
  \item \textbf{Scope.} The system handles English prose of at least 100
        words. It does not localise AI-written passages inside mixed
        documents, does not identify which model wrote a text, and has not been
        evaluated on code, poetry written for assignments, or very short
        answers.
  \item \textbf{Paraphrasing and editing.} The most accurate system on
        paraphrased MAGE text (the thorough cascade) detects only three in four
        paraphrased AI texts; a determined user who paraphrases or rewrites the
        output can evade detection.
  \item \textbf{Transformer brittleness.} The second stage was fine-tuned only
        on the three training benchmarks and inherits their genre biases,
        which is the main source of false positives on escalated essays.
\end{itemize}

\textbf{Limitations of the study.}
\begin{itemize}
  \item \textbf{Generators.} The corpus contains public benchmark generations
        up to GPT-4, Llama-2, Mistral and MPT. Custom generations from current
        GPT-4-class, Claude and Llama~3 models were not produced, so the results
        may not carry over to the newest models.
  \item \textbf{Sample sizes.} The held-out Mistral set (744 documents), the
        attack set (690) and above all the fairness test set (128 essays) are
        small; differences of a few percentage points on them are within
        sampling noise.
  \item \textbf{Single run.} Every model was trained once with one random seed;
        variation between training runs is not measured, and the confidence
        intervals reflect only test-set sampling.
  \item \textbf{Fairness proxy.} Native and non-native writing are represented by
        three essay collections that also differ in topic, age group and
        genre, so part of the measured disparity may be due to these
        differences rather than to the writer's first language.
  \item \textbf{Alternative encoder.} DeBERTa-v3 was implemented but not
        trained, so the comparison between encoders was not carried out.
\end{itemize}

% ---------------------------------------------------------------------
\section{Future Scope}
\label{sec:future}
% ---------------------------------------------------------------------

\textbf{A genuinely cheap first stage.} The most important next step follows
directly from the efficiency result. The reference statistics can be computed
with a much smaller language model, for example a distilled GPT-2 or a
70--160 million parameter model with a smaller vocabulary, or with weak
unigram and trigram models as in Ghostbuster~\cite{verma2024ghostbuster}; the
number of scored tokens can be capped at 512; and spaCy POS tagging can be
replaced by a faster tagger or dropped, since the twenty POS bigrams together
contribute under 10\% of the head's gain while spaCy dominates the stage's cost
on the GPU. A two-tier first stage, in which text-only stylometric
features decide the clearest documents before any language model is run, would
reduce the cost further. Each variant can be evaluated with the existing
pipeline by changing one configuration value and repeating the feature
extraction.

\textbf{A more robust second stage.} The transformer's collapse on
out-of-distribution text and on student essays calls for training data from
more genres (school and university essays in particular), adversarial or
paraphrase-augmented training as in RADAR~\cite{hu2023radar}, and evaluation
of DeBERTa-v3. Because the cascade escalates exactly the hard documents, the
second stage could also be trained with extra weight on them.

\textbf{Current generators.} Running the provided generation script for
current GPT-4-class, Claude and Llama~3 models, and repeating the
leave-one-generator-out evaluation for them, would show how well the detector
transfers to the models students actually use today.

\textbf{Better fairness evaluation.} Larger and better-matched native and
non-native essay collections on the same prompts would separate the effect of
the writer's first language from genre effects. Group-aware thresholds or a
per-institution recalibration (already supported by the
\texttt{retrain-heuristic} and \texttt{tune} commands) could then reduce the
remaining disparity.

\textbf{Finer-grained output.} Sentence-level scoring would allow the system to
highlight which passages of a mixed document look machine-generated, which is
both more useful to a reviewer and more informative to the student.

\textbf{Deployment.} Integration with a learning management system, a
reviewer workflow that records decisions and their outcomes, and periodic
recalibration on new labelled data would turn the prototype into a tool that an
institution can operate responsibly.
